\documentclass[11pt]{article}

\usepackage[margin=1in]{geometry}
\usepackage{amsmath, amssymb}
\usepackage{xcolor}
\usepackage{listings}
\usepackage{hyperref}
\usepackage{array}
\usepackage{booktabs}

\hypersetup{
    colorlinks=true,
    linkcolor=blue!60!black,
    urlcolor=blue!60!black,
    citecolor=blue!60!black
}

\lstdefinestyle{bugpython}{
  language=Python,
  basicstyle=\ttfamily\small,
  keywordstyle=\color{blue!60!black},
  stringstyle=\color{green!40!black},
  commentstyle=\color{gray!70!black},
  showstringspaces=false,
  breaklines=true,
  frame=single,
  rulecolor=\color{gray!30},
  columns=fullflexible,
  keepspaces=true
}

\title{SymPy Bug Report}
\date{}

\begin{document}

\author{}
\maketitle

\section{Bug 18: \texttt{sinh} Falsely Evaluates a Nonzero Split Imaginary-Period Argument to Zero}

\subsection{Minimal Reproducer}

\medskip

\begin{lstlisting}[style=bugpython]
from sympy import I, pi, simplify, sinh

expr = sinh(I*pi + I*(pi - 1))
combined = sinh(I*(2*pi - 1))

print("sinh(I*pi + I*(pi - 1)) =", expr)
print("sinh(I*(2*pi - 1)) =", combined)
print("difference =", simplify(expr - combined))
\end{lstlisting}

\subsection{Observed Behavior}

\medskip

\begin{lstlisting}[style=bugpython]
sinh(I*pi + I*(pi - 1)) = 0
sinh(I*(2*pi - 1)) = -I*sin(1)
difference = I*sin(1)
cmath expected = -0.8414709848078966j
\end{lstlisting}

SymPy evaluates the split argument to zero even though the algebraically combined
argument evaluates to \(-I\sin(1)\). The recorded numerical oracle is also
nonzero, so this is not merely a different but equivalent symbolic form.

\subsection{Expected Behavior}

The expression \(\sinh(I\pi + I(\pi - 1))\) should evaluate to
\(-I\sin(1)\), with numerical value approximately
\(-0.8414709848078965\,I\). Equivalently, it should agree with
\(\sinh(I(2\pi - 1))\).

\subsection{Mathematical Explanation}

The two arguments are identical:
\[
  I\pi + I(\pi - 1) = I(2\pi - 1).
\]
For real \(t\), the identity \(\sinh(It)=I\sin(t)\) gives
\[
  \sinh(I(2\pi - 1)) = I\sin(2\pi - 1) = -I\sin(1).
\]
Since \(\sin(1) \ne 0\), the correct value is nonzero. The function
\(\sinh\) is entire, so there is no branch convention or missing integration
constant that could make the split expression evaluate to zero while the
combined expression evaluates to \(-I\sin(1)\).

\subsection{Source-Level Diagnosis}

The diagnosis identifies the relevant source locations as follows.
\begin{center}
\renewcommand{\arraystretch}{1.15}
\begin{tabular}{@{}>{\raggedright\arraybackslash}p{0.58\linewidth}>{\raggedright\arraybackslash}p{0.32\linewidth}@{}}
\toprule
\textbf{File} & \textbf{Relevant function or helper} \\
\midrule
\nolinkurl{sympy/functions/elementary/hyperbolic.py} & \texttt{sinh.eval}, lines 189--208 \\
\nolinkurl{sympy/functions/elementary/trigonometric.py} & \texttt{sin.eval}, lines 301--382 \\
\nolinkurl{sympy/functions/elementary/trigonometric.py} & \texttt{\_peeloff\_pi}, lines 154--159 \\
\bottomrule
\end{tabular}
\end{center}

The traced call path starts in \texttt{sinh.eval}. SymPy extracts the
coefficient of \(I\) with \nolinkurl{_imaginary_unit_as_coefficient} and
rewrites the call as \(I\sin(\texttt{i\_coeff})\). For the unevaluated split
argument, that coefficient is the addend structure \(({-1}+\pi)+\pi\). Then
\texttt{sin.eval} calls \texttt{\_peeloff\_pi}.

The source-level issue is in \texttt{\_peeloff\_pi}. The helper iterates over
\texttt{Add.make\_args(arg)} and uses \texttt{a.coeff(pi)} to remove rational
multiples of \(\pi\). For the nested addend \(-1+\pi\),
\texttt{a.coeff(pi)} is \(1\). The diagnosis reports that this increments the
peeled \(\pi\)-coefficient but fails to preserve the non-\(\pi\) remainder
\(-1\) in the residual terms.

As a result, the argument \(\pi+(\pi-1)\) is treated as the exact multiple
\(2\pi\). The downstream trigonometric simplification evaluates
\(\sin(2\pi)\) to zero, and the hyperbolic rewrite propagates this as
\(\sinh(I\pi+I(\pi-1)) = I\cdot 0 = 0\). A plausible fix direction supported
by the diagnosis is to peel a \(\pi\) multiple from an addend only when the
whole addend is a pure rational multiple of \(\pi\), or otherwise subtract the
peeled multiple and keep any nonzero remainder in the residual terms.

\subsection{Additional Failing Instances}

The independent verification checks the family
\[
  \sinh(I\pi + I(\pi-k)), \qquad k \in \{1,2,3,4,5,6\}.
\]
For every listed integer offset, the same identity applies:
\[
  \sinh(I\pi + I(\pi-k))=\sinh(I(2\pi-k))
  = I\sin(2\pi-k) = -I\sin(k).
\]
The expected value is therefore nonzero for these integer offsets, but the
dropped-remainder mechanism makes the split argument look like an exact
\(2\pi\) trigonometric argument.

\begin{center}
\renewcommand{\arraystretch}{1.3}
\begin{tabular}{@{}>{\raggedright\arraybackslash}p{0.44\linewidth}>{\raggedright\arraybackslash}p{0.23\linewidth}>{\raggedright\arraybackslash}p{0.23\linewidth}@{}}
\toprule
\textbf{Input / Expression} & \textbf{SymPy behavior} & \textbf{Expected behavior} \\
\midrule
\(\sinh(I\pi+I(\pi-1))\) & \(0\) & \(-I\sin(1)\) \\
\(\sinh(I\pi+I(\pi-2))\) & \(0\) & \(-I\sin(2)\) \\
\(\sinh(I\pi+I(\pi-3))\) & \(0\) & \(-I\sin(3)\) \\
\(\sinh(I\pi+I(\pi-4))\) & \(0\) & \(-I\sin(4)\) \\
\(\sinh(I\pi+I(\pi-5))\) & \(0\) & \(-I\sin(5)\) \\
\(\sinh(I\pi+I(\pi-6))\) & \(0\) & \(-I\sin(6)\) \\
\bottomrule
\end{tabular}
\end{center}

\subsection{Related Issues Assessment}

The bug-finding harness performed a best-effort deduplication check. The closest
related public issues involve trigonometric periodicity and the same
\texttt{\_peeloff\_pi} helper: SymPy issue \#12328 discusses
\texttt{\_peeloff\_pi} in a assumptions-metadata case involving
\(\sin(m\pi+\pi)\), while issue \#17976 concerns incomplete periodic
simplification such as \(\sin(2n\pi+4)\). These are related at the subsystem
level but do not report the exact split imaginary-period reproducer, the wrong
zero result, or the dropped non-\(\pi\) remainder from \(\pi-1\). The
deduplication verdict is therefore a new instance of a known failure mode, and
the concrete reproducer remains individually actionable as a SymPy issue and
regression test.

\subsection{Proposed SymPy Regression Test}

\medskip

\begin{lstlisting}[style=bugpython]
import pytest

from sympy import I, pi, simplify, sin, sinh


@pytest.mark.parametrize("offset", [1, 2, 3, 4, 5, 6])
def test_sinh_split_imaginary_period_preserves_remainder(offset):
    assert simplify(sinh(I*pi + I*(pi - offset)) + I*sin(offset)) == 0
\end{lstlisting}

\end{document}
